% =====================================================================
%  CHAPTER 11: 食源性疾病及其预防
% =====================================================================
\chapter{食源性疾病及其预防}
\setcounter{chapter}{11}
\chapterintro{
\keyword{掌握}食源性疾病的定义；掌握食物中毒的定义、特点；掌握细菌性食物中毒的发病原因、特点及常见类型（沙门氏菌、金黄色葡萄球菌、副溶血性弧菌、肉毒梭菌）；掌握河豚鱼中毒、毒蕈中毒；掌握亚硝酸盐食物中毒、有机磷农药中毒。\keyword{熟悉}食物过敏。
}

% =====================================================================
%  第一节：食源性疾病
% =====================================================================
\section{食源性疾病}

\begin{keybox}
\textbf{食源性疾病（foodborne disease）的定义：}
食品中致病因素进入人体引起的感染性、中毒性疾病。
\end{keybox}

食源性疾病包括食物中毒、食源性肠道传染病、食源性寄生虫病、人畜共患传染病、食物过敏、食物污染物引起的慢性中毒性疾病等。其中，食物中毒是最常见的一类。

% =====================================================================
%  第二节：食物中毒
% =====================================================================
\section{食物中毒}

\subsection{食物中毒的定义}

\begin{keybox}
\textbf{食物中毒（food poisoning）的定义：}
食用被有毒有害物质污染的食品或食用含有毒有害物质的食品后出现的急性、亚急性疾病。
\end{keybox}

\subsection{食物中毒的特点}

\begin{keybox}
\textbf{食物中毒的四大临床特点（掌握）：}
\begin{enumerate}[leftmargin=*]
    \item \textbf{发病与特定食物有关：}所有发病者均在近期内进食过同一种可疑食物，未进食该食物者不发病。停止食用该食物后，发病立即停止。
    \item \textbf{潜伏期短，发病曲线单峰型：}短时间内多人同时或相继发病，发病急剧，呈暴发性。发病曲线呈突然上升又迅速下降的单峰型，无传染病流行的余波（拖尾现象）。
    \item \textbf{临床表现相似：}以急性胃肠炎症状为主，最常见为消化道症状（腹痛、腹泻、恶心、呕吐）。
    \item \textbf{一般没有人与人间的直接传染：}停止进食致病食品后，发病立即下降，无二代病例出现。
\end{enumerate}
\end{keybox}

\subsection{食物中毒的分类}

食物中毒按致病因素可分为以下五类：
\begin{enumerate}[leftmargin=*]
    \item \textbf{细菌性食物中毒：}发病率最高、中毒人数最多
    \item \textbf{真菌毒素和霉变食品中毒}
    \item \textbf{动物性食物中毒：}如河豚鱼中毒
    \item \textbf{有毒植物食物中毒：}如毒蕈中毒
    \item \textbf{化学性食物中毒：}如亚硝酸盐中毒、有机磷农药中毒
\end{enumerate}

% =====================================================================
%  第三节：细菌性食物中毒
% =====================================================================
\section{细菌性食物中毒}

\begin{defbox}
\textbf{细菌性食物中毒：}摄入被致病菌或其毒素污染的食品后发生的急性或亚急性疾病。是最常见的食物中毒。
\end{defbox}

\subsection{发病原因}

\begin{keybox}
\textbf{细菌性食物中毒的三大发病原因（掌握）：}
\begin{enumerate}[leftmargin=*]
    \item \textbf{致病菌污染食品：}原料污染（生前感染或宰后污染）、加工过程污染、从业人员带菌污染。
    \item \textbf{食品的贮存方式不当：}在较高温度下存放，使致病菌大量生长繁殖和/或产生毒素。
    \item \textbf{食用前未彻底加热处理：}或生熟交叉污染，使食品中心温度未达到灭菌要求。
\end{enumerate}
\end{keybox}

\subsection{流行病学特点}

\begin{keybox}
\textbf{细菌性食物中毒的特点（掌握）：}
\begin{enumerate}[leftmargin=*]
    \item \textbf{潜伏期较短：}摄入致病菌或其毒素后，短时间内即可发病。
    \item \textbf{多呈暴发性：}短时间内多人同时或相继发病。
    \item \textbf{发病季节性明显：}夏秋季（5～10月）高发，此时气温高、湿度大，细菌易繁殖。
    \item \textbf{以胃肠道症状为主：}腹痛、腹泻、恶心、呕吐等急性胃肠炎表现。
    \item \textbf{病死率较低：}\keyword{肉毒梭菌中毒除外！}一般细菌性食物中毒如能及时治疗，预后良好。
\end{enumerate}
\end{keybox}

\subsection{常见类型}

\subsubsection{一、沙门氏菌食物中毒}

\begin{keybox}
\textbf{沙门氏菌食物中毒（掌握）：}
\begin{itemize}[leftmargin=*]
    \item \textbf{来源：}主要为动物性食品（畜禽肉、蛋类、乳类）。其中畜肉类及其制品居首位，其次为禽肉、蛋类、乳类。
    \item \textbf{病原特点：}沙门菌属（Salmonella），革兰阴性杆菌。\keyword{不分解蛋白质、不产生靛基质}（吲哚），因此\textbf{污染食品后无感官性状改变}——这是其最危险的特性。
    \item \textbf{中毒机制：}\keyword{活菌 + 内毒素}（感染型）。活菌随食物进入肠道后大量繁殖，侵入肠黏膜上皮细胞，引起炎症反应并释放内毒素。
    \item \textbf{临床特点：}\keyword{潜伏期12～36h}（一般12～14h）。\textbf{发热}、恶心呕吐、\textbf{腹痛腹泻（水样便或黄绿色稀便）}。
    \item \textbf{高发季节：}夏秋季（5～10月）占全年的80\%。
    \item \textbf{灭菌要求：}肉块中心温度$\ge$80$^\circ$C持续12分钟；鸡蛋煮沸8～10分钟。
\end{itemize}
\end{keybox}

\subsubsection{二、金黄色葡萄球菌食物中毒}

\begin{keybox}
\textbf{金黄色葡萄球菌食物中毒（掌握）：}
\begin{itemize}[leftmargin=*]
    \item \textbf{来源：}化脓性皮肤病或上呼吸道感染患者\textbf{对食品的污染}。人体鼻咽部带菌率20\%～50\%，皮肤化脓性病灶含大量致病性葡萄球菌。\keyword{带菌者污染食品是最重要的污染来源！}
    \item \textbf{中毒机制：}\keyword{肠毒素（毒素型）}。金黄色葡萄球菌在食品中繁殖时产生肠毒素，该肠毒素\textbf{耐热}——\keyword{100$^\circ$C加热2h才被破坏}，普通烹调方法难以彻底破坏。
    \item \textbf{临床特点：}\keyword{潜伏期1～6h}（极短，最短仅1h）。\textbf{剧烈反复呕吐}（最突出的症状！）、\textbf{上腹剧痛}、水样腹泻。\keyword{体温一般正常}。
    \item \textbf{预防：}\keyword{有局部化脓性感染或上呼吸道感染者应暂时调离食品工作岗位！}
\end{itemize}
\end{keybox}

\subsubsection{三、副溶血性弧菌食物中毒}

\begin{keybox}
\textbf{副溶血性弧菌食物中毒（掌握）：}
\begin{itemize}[leftmargin=*]
    \item \textbf{来源：}主要为\textbf{海产品}（鱼类、贝类、虾蟹等）。墨鱼的带菌率可达93\%。7～9月为高发季节。\keyword{嗜盐菌}——最适生长盐浓度3.5\% NaCl，淡水环境中迅速死亡。
    \item \textbf{中毒机制：}\keyword{活菌 + 溶血毒素}（混合型）。神奈川试验（Kanagawa Test）阳性——90\%的致病性菌株为神奈川现象阳性。
    \item \textbf{临床特点：}\keyword{潜伏期6～12h}。\textbf{上腹阵发性绞痛}（最突出的特点！）、\textbf{水样便或洗肉水样便}（血水样便）、\textbf{里急后重不明显}。发热较明显（37.5～39.5$^\circ$C）。
    \item \textbf{杀菌条件：}56$^\circ$C加热5分钟或90$^\circ$C加热1分钟即可杀灭；\keyword{含1\%食醋处理5分钟即可杀灭}——食用海鲜蘸醋对预防有意义。
\end{itemize}
\end{keybox}

\begin{exambox}
\textbf{副溶血性弧菌三大特征：}嗜盐性（3.5\% NaCl最适）、神奈川现象（90\%致病株阳性）、血水样便（洗肉水样便）。
\end{exambox}

\subsubsection{四、肉毒梭菌食物中毒}

\begin{keybox}
\textbf{肉毒梭菌食物中毒（掌握）}\difficulty：
\begin{itemize}[leftmargin=*]
    \item \textbf{来源：}家庭自制发酵食品（豆酱、臭豆腐、豆豉等）和罐头食品。中国以\textbf{植物性食品为主}（占91.48\%），新疆地区发病率最高。
    \item \textbf{中毒机制：}\keyword{外毒素——肉毒毒素（botulinum toxin）}，是\textbf{已知最剧烈的神经毒素}。毒素经小肠激活后，作用于神经肌肉接头，阻止乙酰胆碱释放，导致肌肉麻痹。人的致死剂量约10$^{-9}$mg/kg体重，毒性比KCN强10\,000倍。
    \item \textbf{临床特点：}\keyword{潜伏期12～48h}。以\textbf{运动神经麻痹症状为主}，胃肠道症状极少见。主要表现：\textbf{对称性颅神经损害}——\textbf{视力模糊、复视、眼睑下垂、吞咽困难、语言障碍、呼吸困难}（致死主要原因）。\textbf{病死率极高}。
    \item \textbf{芽胞耐热性：}\keyword{干热180$^\circ$C 5～15分钟，湿热100$^\circ$C 5小时，高压蒸汽灭菌121$^\circ$C 30分钟。}
    \item \textbf{治疗：}\keyword{尽早使用多价肉毒抗毒素血清}（A、B、E型），同时对症支持治疗。
    \item \textbf{婴儿肉毒中毒：}\keyword{蜂蜜是1岁以下婴儿肉毒中毒的嫌疑食品}，主要为E型。
    \item \textbf{预防：}家庭自制密封发酵食品为高危风险食品；罐头食品鼓盖/漏气绝对不要购买食用。
\end{itemize}
\end{keybox}

\begin{exambox}
\textbf{肉毒梭菌与其他细菌性食物中毒的根本区别：}胃肠道症状极少见，不以呕吐腹泻为主，主要表现为运动神经麻痹（对称性颅神经损害）。芽胞耐热性极强——湿热100$^\circ$C需5小时！中国以发酵豆制品为主要中毒食品，新疆高发！
\end{exambox}

\begin{tipbox}
\textbf{四种细菌性食物中毒对比速记表：}

\bgroup
\renewcommand{\arraystretch}{1.3}
\begin{tabularx}{\textwidth}{lXXXX}
\toprule
\textbf{病原体} & \textbf{来源} & \textbf{中毒机制} & \textbf{潜伏期} & \textbf{核心特征} \\
\midrule
沙门氏菌 & 动物性食品（肉、蛋、乳） & 活菌+内毒素（感染型） & 12～36h & 发热+水样便；食物无感官改变 \\
\addlinespace
金葡菌 & 带菌者（化脓/呼吸道感染） & 肠毒素（毒素型），耐热 & \textbf{1～6h} & 剧烈呕吐最突出；体温正常 \\
\addlinespace
副溶血性弧菌 & 海产品 & 活菌+溶血毒素（混合型） & 6～12h & 腹绞痛+洗肉水样便；嗜盐 \\
\addlinespace
肉毒梭菌 & 家庭发酵食品、罐头 & 肉毒毒素（神经毒素） & 12～48h & \textbf{运动神经麻痹}；病死率极高 \\
\bottomrule
\end{tabularx}
\egroup
\end{tipbox}

% =====================================================================
%  第四节：有毒动植物中毒
% =====================================================================
\section{有毒动植物中毒}

\subsection{河豚鱼中毒}

\begin{keybox}
\textbf{河豚鱼中毒（掌握）：}
\begin{itemize}[leftmargin=*]
    \item \textbf{有毒成分：}\keyword{河豚毒素（Tetrodotoxin, TTX）}，非蛋白质神经毒素。\textbf{耐热！煮沸、盐腌、日晒均不能破坏。}毒素分布：\keyword{卵巢毒性最强}，肝脏次之，血液、眼球、皮肤亦含毒素。新鲜洗净的肌肉基本无毒。每年\textbf{春季（2～5月）}为卵巢发育期，毒性最强。
    \item \textbf{中毒机制：}TTX\textbf{阻断Na$^+$通道}，阻止神经细胞膜钠离子内流，使神经传导阻滞。先麻痹神经末梢和感觉神经，后麻痹运动神经和中枢神经系统。
    \item \textbf{临床特点：}\keyword{潜伏期10min～3h}。发病急而剧烈，按顺序出现：
    \begin{enumerate}[leftmargin=*]
        \item \textbf{先感觉神经麻痹：}手指、\textbf{口唇、舌尖发麻}刺痛
        \item \textbf{后运动神经麻痹：}恶心、呕吐、四肢无力、肌肉麻痹、全身瘫痪
        \item \textbf{最后呼吸麻痹死亡：}呼吸衰竭是致死主要原因。通常死于发病后4～6小时，最快1.5小时。
    \end{enumerate}
    \item \textbf{治疗：}无特效解毒药！尽早洗胃 + 导泻 + 对症支持治疗（维持呼吸、升压等）。
    \item \textbf{预防：}\keyword{禁止销售河豚鱼。}严禁餐饮单位经营加工河豚鱼，加强宣传教育。
\end{itemize}
\end{keybox}

\subsection{毒蕈中毒}

\begin{keybox}
\textbf{毒蕈中毒（掌握）}\difficulty：
\begin{itemize}[leftmargin=*]
    \item \textbf{基本概况：}我国已知毒蕈约80余种，含剧毒可致死的不超过10种。不同类型毒素有\textbf{不同的临床表现}。
    \item \textbf{临床分型：}
    \begin{enumerate}[leftmargin=*]
        \item \textbf{胃肠炎型：}潜伏期短（10min～6h），恶心、呕吐、腹痛、腹泻。对症支持治疗，预后良好。
        \item \textbf{神经精神型：}毒蝇碱、裸盖菇素等。出汗、流涎、瞳孔缩小、幻视、精神错乱。\keyword{阿托品（Atropine）}拮抗治疗。
        \item \textbf{溶血型：}鹿花蕈毒素。溶血性贫血、黄疸、血红蛋白尿。\keyword{肾上腺皮质激素}治疗。
        \item \textbf{肝肾损害型（最严重！）：}毒肽（Phallotoxins）和毒伞肽（Amanitins）引起。\textbf{病死率最高。}潜伏期长（10～24h），病情发展可分为6期：
        \begin{enumerate}[leftmargin=*]
            \item 潜伏期（10～24h）
            \item 胃肠炎期
            \item \textbf{假愈期}（胃肠炎症状缓解但肝脏损害已开始——最易被忽视！）
            \item 内脏损害期（肝、肾、脑、心损害）
            \item 精神症状期（肝性昏迷、中毒性脑病）
            \item 恢复期
        \end{enumerate}
        抢救：\keyword{巯基解毒剂（二巯基丙磺酸钠、二巯基丁二酸钠等）是抢救关键！}
        \item 类光过敏型：类似日光性皮炎，身体暴露部位明显肿胀、疼痛，嘴唇肿胀外翻。
    \end{enumerate}
    \item \textbf{鹅膏菌中毒：}食用后\textbf{10小时以内彻底洗胃}是关键！
\end{itemize}
\end{keybox}

\begin{exambox}
\textbf{肝肾损害型毒蕈中毒的"假愈期"是重要考点：}胃肠炎症状暂时缓解后病情突然恶化，患者和家属误以为已康复，导致错过抢救时机。巯基解毒剂是关键抢救药物。
\end{exambox}

% =====================================================================
%  第五节：化学性食物中毒
% =====================================================================
\section{化学性食物中毒}

\subsection{亚硝酸盐食物中毒}

\begin{keybox}
\textbf{亚硝酸盐食物中毒（掌握）：}
\begin{itemize}[leftmargin=*]
    \item \textbf{中毒原因：}
    \begin{enumerate}[leftmargin=*]
        \item 误将亚硝酸盐当作食盐使用（最危险！建筑工地常见）
        \item 进食大量含硝酸盐/亚硝酸盐较高的蔬菜——新鲜腌制的蔬菜（"暴腌菜"），\keyword{腌制第4～8天亚硝酸盐含量达高峰}
        \item 肠源性青紫症（肠道功能紊乱时，肠道细菌将硝酸盐还原为亚硝酸盐）
        \item 某些井水含硝酸盐过高（苦井水）
    \end{enumerate}
    \item \textbf{中毒机制：}亚硝酸盐将\textbf{血红蛋白（Hb）中的Fe$^{2+}$氧化为Fe$^{3+}$}，形成\textbf{高铁血红蛋白（MetHb）}，失去携氧能力，导致组织缺氧，出现\textbf{肠源性青紫症（发绀）}。
    \item \textbf{临床特点：}\keyword{潜伏期10min～3h}。\textbf{口唇、指尖发绀}（肠源性青紫症），头晕、乏力、恶心、呕吐。严重者呼吸困难、昏迷、循环衰竭死亡。
    \item \textbf{特效解毒剂：}\keyword{美蓝（亚甲蓝，Methylene Blue）} + \textbf{维生素C}。小剂量美蓝（1～2 mg/kg）将MetHb还原为Hb；维生素C作为辅助还原剂。
    \item \textbf{安全食用腌菜：}\keyword{腌制20天以后亚硝酸盐基本消失}，方可安全食用。
\end{itemize}
\end{keybox}

\begin{tipbox}
\textbf{暴腌菜风险：}腌制第4～8天亚硝酸盐含量达峰值，此时食用中毒风险最大！腌制20天以上才安全。误将亚硝酸盐当食盐是化学性食物中毒的最常见原因。
\end{tipbox}

\subsection{有机磷农药中毒}

\begin{keybox}
\textbf{有机磷农药中毒（掌握）：}
\begin{itemize}[leftmargin=*]
    \item \textbf{中毒原因：}误食农药拌过的种子；误把有机磷农药当作食用油或酱油食用；食用喷洒农药后未经安全间隔期即采摘的蔬菜水果；误食被农药毒杀的家禽家畜。
    \item \textbf{中毒机制：}\keyword{抑制胆碱酯酶（ChE）活性}，使其失去催化水解乙酰胆碱（ACh）的能力 $\to$ 大量ACh在体内蓄积 $\to$ 胆碱能神经过度兴奋 $\to$ 出现毒蕈碱样症状、烟碱样症状和中枢神经系统症状。
    \item \textbf{临床特点：}\keyword{潜伏期30min～3h}。
    \begin{enumerate}[leftmargin=*]
        \item \textbf{毒蕈碱样症状（M样症状）：}\keyword{瞳孔缩小}、\keyword{流涎}、\keyword{大汗}、恶心呕吐、腹痛腹泻、支气管痉挛、呼吸困难。
        \item \textbf{烟碱样症状（N样症状）：}\keyword{肌束震颤}（从眼睑、面部小肌群开始，逐渐发展到全身）、肌肉无力甚至瘫痪。
        \item \textbf{中枢神经系统症状：}头晕、头痛、烦躁不安、抽搐、昏迷。严重者因呼吸中枢抑制和呼吸肌麻痹导致呼吸衰竭死亡。
    \end{enumerate}
    \item \textbf{全血胆碱酯酶活性分级：}
    \begin{itemize}[leftmargin=*]
        \item 轻度中毒：ChE活性降至正常值的50\%～70\%
        \item 中度中毒：ChE活性降至正常值的30\%～50\%
        \item 重度中毒：ChE活性降至正常值的30\%以下
    \end{itemize}
    \item \textbf{特效解毒剂：}\keyword{阿托品（Atropine） + 氯解磷定（Pralidoxime Chloride）/碘解磷定}。阿托品拮抗毒蕈碱样症状（对抗ACh蓄积），解磷定使被抑制的胆碱酯酶复活。
    \item \textbf{预防：}严格执行农药安全使用规定，遵守安全间隔期；加强农药保管；严禁食用被农药毒死的动物。
\end{itemize}
\end{keybox}

\begin{tipbox}
\textbf{有机磷中毒症状速记：}M样 = 毒蕈碱样 —— 流涎、大汗、瞳孔缩小（"水"的症状）；N样 = 烟碱样 —— 肌束震颤（肌肉症状）。特效解毒剂：阿托品（拮抗M样症状）+ 解磷定（复活ChE）。
\end{tipbox}

% =====================================================================
%  第六节：食物过敏
% =====================================================================
\section{食物过敏}

\begin{defbox}
\textbf{食物过敏（food allergy）的定义：}
机体免疫系统对进入体内的食物中的某些非感染性物质（过敏原）产生的异常免疫应答，导致生理功能紊乱和/或组织损伤。
\end{defbox}

\subsection{常见致敏食物}

\textbf{八大类主要致敏食物：}
\begin{enumerate}[leftmargin=*]
    \item 牛奶（Milk）
    \item 鸡蛋（Eggs）
    \item 花生（Peanuts）
    \item 大豆（Soy）
    \item 谷物（小麦等，Grains/Wheat）
    \item 鱼类（Fish）
    \item 甲壳类（Crustaceans）
    \item 坚果类（Tree Nuts）
\end{enumerate}

\subsection{食物过敏的特点}

\begin{enumerate}[leftmargin=*]
    \item 任何食物均可诱发过敏，但少数食物是主要过敏原
    \item 仅有部分食物成分具有变应原性，特定蛋白质为主要致敏成分
    \item 不同食物间可发生交叉过敏反应
    \item 食物加工可改变致敏性：加热可使部分过敏原变性降低致敏性，但某些热稳定型过敏原（如牛奶中$\beta$-乳球蛋白）不受加热影响
\end{enumerate}

\subsection{卫生假说（Hygiene Hypothesis）}

\begin{tipbox}
\textbf{卫生假说：}免疫系统在进化过程中为应对肠道寄生虫而设计，现代社会中缺乏肠道寄生虫感染，导致免疫系统过度活跃，对无害食物抗原产生异常免疫应答。这解释了工业化国家过敏性疾病高发的原因。
\end{tipbox}

% =====================================================================
%  复习题
% =====================================================================
\reviewcontent{
\section{复习题}

\subsection*{一、名词解释}

\begin{enumerate}[leftmargin=*, itemsep=6pt]
    \item 食源性疾病（foodborne disease）
    \item 食物过敏（food allergy）
    \item 食物中毒（food poisoning）
    \item 细菌性食物中毒
    \item 肠源性青紫症
    \item 高铁血红蛋白血症（Methemoglobinemia）
    \item 毒蕈中毒假愈期
\end{enumerate}

\subsection*{二、简答题}

\begin{enumerate}[leftmargin=*, itemsep=8pt]
    \item 简述食源性疾病和食物中毒的定义。
    \item 简述食物中毒的四大临床特点。
    \item 简述细菌性食物中毒的三大发病原因和五大特点。
    \item 比较沙门氏菌、金黄色葡萄球菌、副溶血性弧菌、肉毒梭菌四种食物中毒的来源、中毒机制、潜伏期和临床特点。
    \item 金黄色葡萄球菌肠毒素的耐热性如何？为什么带菌者管理对预防金葡菌食物中毒至关重要？
    \item 简述肉毒梭菌毒素的作用机制，并说明为什么肉毒梭菌中毒的病死率极高。
    \item 简述河豚毒素（TTX）的中毒机制和河豚鱼中毒的临床发展过程。
    \item 简述毒蕈中毒的临床分型，以及肝肾损害型的"假愈期"有何临床意义。
    \item 简述亚硝酸盐食物中毒的中毒机制、临床特点和特效解毒剂。
    \item 简述有机磷农药中毒的中毒机制和三大类临床表现。
    \item 简述有机磷农药中毒时全血胆碱酯酶活性的三级变化及相应的特效解毒剂。
    \item 简述食物过敏的八大类致敏食物及食物过敏的基本特点。
\end{enumerate}

\subsection*{三、案例分析}

\begin{enumerate}[leftmargin=*]
    \item \textbf{【案例分析1】}某工地食堂午餐后约2～4小时，50名工人陆续出现剧烈呕吐、上腹剧痛、水样腹泻，但体温均正常。调查发现午餐食用了凉拌熟肉。请回答：
    \begin{enumerate}[leftmargin=*]
        \item 最可能的中毒病原体是什么？
        \item 该病原体的中毒机制是什么？
        \item 肠毒素的耐热特点如何？
        \item 预防措施中最关键的一点是什么？
    \end{enumerate}

    \item \textbf{【案例分析2】}某家庭聚餐后次日，10人中8人出现视力模糊、复视、眼睑下垂、吞咽困难、发音不清，均无发热和明显胃肠道症状。调查发现食用了家庭自制的密封发酵豆酱。请回答：
    \begin{enumerate}[leftmargin=*]
        \item 最可能的诊断是什么？
        \item 该病原体产生的是何种毒素？其毒力如何？
        \item 应如何治疗？
        \item 如何开展预防？
    \end{enumerate}

    \item \textbf{【案例分析3】}某食堂午餐后1小时左右，多名就餐者出现口唇、指尖发绀，伴头晕、乏力、恶心。调查发现食用了腌制仅5天的暴腌菜和苦井水煮的饭菜。请回答：
    \begin{enumerate}[leftmargin=*]
        \item 最可能的中毒类型是什么？
        \item 其中毒机制是什么？
        \item 特效解毒剂是什么？
        \item 暴腌菜腌制多少天后亚硝酸盐基本消失？
    \end{enumerate}
\end{enumerate}

\subsection*{四、参考答案}

\subsubsection*{一、名词解释参考答案}
\begin{enumerate}[leftmargin=*, itemsep=4pt]
    \item \textbf{食源性疾病（foodborne disease）：}食品中致病因素进入人体引起的感染性、中毒性疾病。
    \item \textbf{食物过敏（food allergy）：}机体免疫系统对进入体内的食物中的某些非感染性物质（过敏原）产生的异常免疫应答，导致生理功能紊乱和/或组织损伤。
    \item \textbf{食物中毒（food poisoning）：}食用被有毒有害物质污染的食品或食用含有毒有害物质的食品后出现的急性、亚急性疾病。
    \item \textbf{细菌性食物中毒：}摄入被致病菌或其毒素污染的食品后发生的急性或亚急性疾病，是最常见的食物中毒。
    \item \textbf{肠源性青紫症：}亚硝酸盐将血红蛋白中的Fe$^{2+}$氧化为Fe$^{3+}$形成高铁血红蛋白，失去携氧能力，口唇、指尖出现发绀，因与肠道细菌将硝酸盐还原为亚硝酸盐有关，故称肠源性青紫症。
    \item \textbf{高铁血红蛋白血症（Methemoglobinemia）：}血红蛋白中的Fe$^{2+}$被氧化为Fe$^{3+}$形成高铁血红蛋白（MetHb），失去携氧能力导致的病理状态。
    \item \textbf{毒蕈中毒假愈期：}肝肾损害型毒蕈中毒病程中的特殊阶段，胃肠炎症状暂时缓解或消失，但肝脏损害已开始进行性发展，最容易被忽视而导致病情恶化。
\end{enumerate}

\subsubsection*{二、简答题参考答案}
\begin{enumerate}[leftmargin=*, itemsep=4pt]
    \item \textbf{食源性疾病和食物中毒的定义：}
    \begin{itemize}[leftmargin=*]
        \item 食源性疾病：食品中致病因素进入人体引起的感染性、中毒性疾病。
        \item 食物中毒：食用被有毒有害物质污染的食品或食用含有毒有害物质的食品后出现的急性、亚急性疾病。
    \end{itemize}

    \item \textbf{食物中毒的四大临床特点：}
    \begin{enumerate}[leftmargin=*]
        \item 发病与特定食物有关
        \item 潜伏期短，发病曲线单峰型
        \item 临床表现相似：最常见为消化道症状（腹痛、腹泻、恶心、呕吐）
        \item 一般没有人与人间的直接传染
    \end{enumerate}

    \item \textbf{细菌性食物中毒的发病原因和特点：}

    三大发病原因：①致病菌污染食品；②食品的贮存方式不当；③食用前未彻底加热处理。

    五大特点：①潜伏期较短；②多呈暴发性；③发病季节性明显（夏秋季5～10月高发）；④以胃肠道症状为主（腹痛、腹泻、恶心、呕吐）；⑤病死率较低（肉毒梭菌除外）。

    \item \textbf{四种细菌性食物中毒比较：}
    \begin{itemize}[leftmargin=*]
        \item \textbf{沙门氏菌：}来源为动物性食品（畜禽肉、蛋类、乳类）；中毒机制为活菌+内毒素（感染型）；潜伏期12～36h；临床特点为发热+水样便/黄绿色稀便，污染食品后无感官改变。
        \item \textbf{金黄色葡萄球菌：}来源为化脓性皮肤病或上呼吸道感染患者的污染；中毒机制为肠毒素（毒素型，耐热，100$^\circ$C 2h破坏）；潜伏期1～6h；临床特点为剧烈反复呕吐+上腹剧痛+水样腹泻，体温一般正常。
        \item \textbf{副溶血性弧菌：}来源为海产品（嗜盐菌）；中毒机制为活菌+溶血毒素（混合型）；潜伏期6～12h；临床特点为上腹阵发性绞痛+水样便或洗肉水样便，里急后重不明显。
        \item \textbf{肉毒梭菌：}来源为家庭自制发酵食品和罐头食品；中毒机制为外毒素（肉毒毒素，已知最剧烈的神经毒素）；潜伏期12～48h；临床特点为对称性颅神经损害（视力模糊、复视、眼睑下垂、吞咽困难、语言障碍、呼吸困难），病死率极高。
    \end{itemize}

    \item \textbf{金葡菌肠毒素耐热性及带菌者管理：}肠毒素耐热，100$^\circ$C加热2h才被破坏，普通烹调方法难以彻底破坏。带菌者（化脓性皮肤病、上呼吸道感染者）是食品污染的最重要来源，因此有局部化脓性感染或上呼吸道感染者应暂时调离食品工作岗位。

    \item \textbf{肉毒梭菌毒素作用机制及高病死率原因：}肉毒毒素作用于神经肌肉接头，阻止乙酰胆碱释放，导致肌肉麻痹。高病死率原因：(1)肉毒毒素是已知最剧烈的神经毒素，致死剂量极低（约10$^{-9}$mg/kg）；(2)主要致死原因为呼吸肌麻痹导致呼吸衰竭；(3)芽胞耐热性极强（湿热100$^\circ$C需5小时），普通加热无法杀灭。

    \item \textbf{河豚毒素（TTX）的中毒机制和临床发展过程：}中毒机制：TTX阻断Na$^+$通道，阻止神经细胞膜钠离子内流，使神经传导阻滞。临床发展过程：先感觉神经麻痹（口唇舌尖发麻）$\to$ 后运动神经麻痹（四肢无力、肌肉麻痹）$\to$ 最后呼吸麻痹死亡。

    \item \textbf{毒蕈中毒的临床分型及假愈期意义：}临床分型：胃肠炎型、神经精神型（阿托品治疗）、溶血型（肾上腺皮质激素）、肝肾损害型（巯基解毒剂，病死率最高）、类光过敏型。假愈期的意义：此期胃肠炎症状暂时缓解，但肝脏损害已开始进行性发展，需继续严密观察和积极治疗，不可因症状缓解而放松警惕。

    \item \textbf{亚硝酸盐食物中毒：}中毒机制：将Hb中的Fe$^{2+}$氧化为Fe$^{3+}$形成MetHb，失去携氧能力。临床特点：潜伏期10min～3h，口唇指尖发绀（肠源性青紫症），头晕乏力。特效解毒剂：美蓝（亚甲蓝）+ 维生素C。

    \item \textbf{有机磷农药中毒机制和临床表现：}中毒机制：抑制胆碱酯酶（ChE）活性，使ACh蓄积，胆碱能神经过度兴奋。三类临床表现：毒蕈碱样症状（瞳孔缩小、流涎、大汗）、烟碱样症状（肌束震颤）、中枢神经系统症状（头晕、昏迷）。

    \item \textbf{有机磷中毒ChE活性分级及特效解毒剂：}轻度：ChE 50\%～70\%；中度：ChE 30\%～50\%；重度：ChE $<$30\%。特效解毒剂：阿托品（拮抗毒蕈碱样症状）+ 氯解磷定/碘解磷定（复活胆碱酯酶）。

    \item \textbf{食物过敏的八大类致敏食物及特点：}八大类：牛奶、鸡蛋、花生、大豆、谷物（小麦等）、鱼类、甲壳类、坚果类。特点：任何食物均可诱发过敏；仅有部分成分具有变应原性；不同食物间可发生交叉过敏反应；食物加工可改变致敏性。
\end{enumerate}

\subsubsection*{三、案例分析参考答案}
\begin{enumerate}[leftmargin=*]
    \item \textbf{【案例分析1】}
    \begin{enumerate}[leftmargin=*]
        \item 最可能的病原体：金黄色葡萄球菌（肠毒素）。
        \item 中毒机制：毒素型——金葡菌在食品中预先产生肠毒素，毒素随食物进入体内引起中毒。
        \item 肠毒素耐热，100$^\circ$C加热2h才被破坏，普通烹调难以彻底破坏。
        \item 最关键预防措施：有局部化脓性感染或上呼吸道感染者应暂时调离食品工作岗位。
    \end{enumerate}

    \item \textbf{【案例分析2】}
    \begin{enumerate}[leftmargin=*]
        \item 最可能的诊断：肉毒梭菌毒素食物中毒。
        \item 肉毒毒素（外毒素），是已知最剧烈的神经毒素。毒性比KCN强10\,000倍，人的致死剂量约10$^{-9}$mg/kg。
        \item 尽早使用多价肉毒抗毒素血清（A、B、E型），同时对症支持治疗。
        \item 预防：家庭自制密封发酵食品为高危风险食品；罐头食品鼓盖/漏气绝对不要食用；蜂蜜不宜喂养1岁以下婴儿。
    \end{enumerate}

    \item \textbf{【案例分析3】}
    \begin{enumerate}[leftmargin=*]
        \item 最可能的诊断：亚硝酸盐食物中毒。
        \item 中毒机制：亚硝酸盐将Hb中的Fe$^{2+}$氧化为Fe$^{3+}$，形成MetHb，失去携氧能力，导致肠源性青紫症。
        \item 特效解毒剂：美蓝（亚甲蓝）+ 维生素C。
        \item 腌制20天以后亚硝酸盐基本消失。
    \end{enumerate}
\end{enumerate}

\newpage
}
